\section{Pointwise inversion throughout the perturbed roof family}
\label{sec:v60-perturbed-density}
We retain the map, probability and roof of
\eqref{eq:v59-markov-baker}--\eqref{eq:v59-continuous-roof}.
Fix $0<\epsilon_0<1/4$. The roof remains positive for
$|\epsilon|\le\epsilon_0$. Define the two original integer counts
\[
 A_m=\sum_{r=1}^m I_r,\qquad
 B_m=\sum_{r=1}^m I_{r-1}I_r,\qquad Q_m=(A_m,B_m).
\]
Write
\begin{equation}\label{eq:v60-joint-data}
 \mu_*=(3/7,2/7),\quad
 \Sigma_*=\frac{1}{343}\begin{pmatrix}204&192\\192&198\end{pmatrix},\quad
 \bar\tau_\epsilon=\frac{24+2\epsilon}{7},\quad
 \sigma_\epsilon^2=\frac{204+384\epsilon+198\epsilon^2}{343}.
\end{equation}
The covariance is positive definite by
\eqref{eq:v59-joint-covariance}. The new proof first keeps both integer
counts and only then projects onto $A_m+\epsilon B_m$. In particular it
does not require a common lattice span for that scalar projection.

\subsection{A two-count local estimate with a positive envelope}
\begin{lemma}[Joint exact-state estimates]
\label{lem:v60-joint-lattice}
For $i,j\in\{0,1\}$, $q\in\Z^2$ and $m\ge1$,
\begin{align}
 \left|m\Pr_i(Q_m=q,I_m=j)
       -\pi_jg_{\Sigma_*}((q-m\mu_*)/\sqrt m)\right|
            &\le C m^{-1/2},                       \label{eq:v60-joint-LLT}\\
 \Pr_i(Q_m=q,I_m=j)
            &\le \frac{C}{m}\exp(-c|q-m\mu_*|^2/m).
                                                        \label{eq:v60-joint-upper}
\end{align}
Consequently, for any fixed $H<\infty$,
\begin{equation}\label{eq:v60-strip-upper}
 \sup_{t,\ |\epsilon|\le\epsilon_0}
 \Pr_i\{|A_m+\epsilon B_m-t|\le H,I_m=j\}
                                \le C_Hm^{-1/2}.
\end{equation}
The stationary initial law satisfies the same estimates after summing
with the original weights $\pi_i$.
\end{lemma}
\begin{proof}
The two-variable Fourier matrix is
$M(z)_{ij}=P_{ij}\exp(i(z_1j+z_2ij))$. If a unit-modulus eigenvalue
exists, equality in the triangle inequality and positivity of all four
transition probabilities force
$e^{i(z_1j+z_2ij)}v_j=\lambda v_i$, with $|v_0|=|v_1|>0$.
The $00$ loop gives $\lambda=1$. The $10$ edge gives $v_0=v_1$,
the $01$ edge gives $e^{iz_1}=1$, and the $11$ loop gives
$e^{iz_2}=1$. Thus the joint lattice is the full $\Z^2$, including
for each specified pair of endpoint states.

Near zero the analytic simple eigenvalue has centered logarithm
$-z^{\mathsf T}\Sigma_*z/2+O(|z|^3)$; its projection entry is
$\pi_j+O(|z|)$. The characteristic polynomial in
\eqref{eq:v59-characteristic-polynomial} supplies the mean and Hessian.
The real part of this logarithm is at most $-c|z|^2$ on a small ball.
Outside it compactness and the preceding phase argument give
exponential decay. In the two-dimensional inverse, the projection
error integrates to $O(m^{-3/2})$; the cubic error is bounded by
$Cm\int |z|^3e^{-cm|z|^2}dz=O(m^{-3/2})$.
Gaussian tails and the complementary power have smaller errors.
The inversion phase has modulus one for every $q$, proving
\eqref{eq:v60-joint-LLT} uniformly in the exact label.

For completeness the upper estimate has no local-limit error. For a
small real tilt $\theta\in\R^2$ use the positive matrix
$M(-i\theta)$. Let $C(\theta)$ be its logarithmic Perron root.
Its Doob transform is a positive stochastic matrix with exactly the
same four edges; its endpoint conjugating factors are uniformly
bounded. The same phase argument, and the positive Hessian of
$C$ on a sufficiently small closed tilt ball, give
\[
 \Pr_i(Q_m=q,I_m=j)
       \le \frac{C}{m}\exp\{mC(\theta)-\theta\cdot q\}.
\]
Indeed the absolute Fourier integral of the tilted matrix powers is
at most $C e^{mC(\theta)}/m$: use the uniform quadratic decrease near
zero and the uniform gap on its compact complement. This is a bound
on a positive coefficient, not subtraction of a signed Gaussian error.
For a feasible count put $v=q-m\mu_*$ and choose $\theta=a v/m$,
where $a>0$ is small and fixed. Feasible counts satisfy $|v|\le C_0m$,
so the tilt lies in the chosen ball. Taylor's inequality
$C(\theta)\le\mu_*\cdot\theta+C_1|\theta|^2$ gives an exponent at
most $-(a-C_1a^2)|v|^2/m$. Infeasible counts have zero probability.
This proves \eqref{eq:v60-joint-upper}.

At fixed $B_m=l$ the strip in \eqref{eq:v60-strip-upper} contains at
most $2H+2$ possible integers $A_m$. Discarding the first coordinate
in the positive Gaussian envelope and summing
$\exp(-c(l-2m/7)^2/m)$ gives $C\sqrt m$. This proves the last bound,
without a rationality or a Diophantine restriction on $\epsilon$.
\end{proof}

On a full word retain the notation $y_m=b_w+a_wy_0$ and $b_m=b_w$;
$0<a_w\le\rho_0^m$, $0\le b_w\le1-a_w$, where $\rho_0=3/4$.
This $b_m$ is the full-word endpoint, not the left endpoint of a suffix
cylinder.

\begin{lemma}[Labelled stable intervals]
\label{lem:v60-joint-cylinders}
For every interval $J$, $q\in\Z^2$ and $1\le L\le m/2$,
\begin{equation}\label{eq:v60-cylinder-LLT}
 \left|m\mu\{Q_m=q,b_m\in J\}
 -g_{\Sigma_*}((q-m\mu_*)/\sqrt m)|J\cap[0,1]|\right|
          \le C\left(\frac{1+L}{\sqrt m}+\rho_0^L\right).
\end{equation}
If in addition $L\le\sqrt m$, the positive bound is
\begin{equation}\label{eq:v60-cylinder-upper}
 \mu\{Q_m=q,b_m\in J\}
 \le \frac{C}{m}e^{-c|q-m\mu_*|^2/m}(|J\cap[0,1]|+\rho_0^L).
\end{equation}
All endpoint conventions for $J$ satisfy these estimates.
\end{lemma}
\begin{proof}
A suffix of $L$ transitions has reward $q_w\in[0,L]^2$, forward
weight $P(w)$, and a terminal stable cylinder of length at most
$\rho_0^L$. If its initial and terminal states are $i,j$, detailed
balance gives $\pi_iP(w)=\pi_j|I_w|$. Its mass at $Q_m=q$ is exactly
\[
 \Pr_\pi(Q_{m-L}=q-q_w,I_{m-L}=i)P(w).
\]
Apply \eqref{eq:v60-joint-LLT} to the prefix. Changing its mean and
scale to those at time $m$ costs $C(1+L)/\sqrt m$ after multiplication
by $m$, uniformly in $q$; the Gaussian and its first derivatives,
including their linear coordinate multiples, are bounded. The errors
sum with $P(w)$, not with the number of cylinders. For all suffixes
ending at either state, $\sum_wP(w)\le2$, since this sum runs over the
two possible starting states. Detailed balance gives the leading
stable length. Inner and outer covers of $J$ require at most two
boundary cylinders per terminal state. This proves
\eqref{eq:v60-cylinder-LLT}.

For the positive statement use \eqref{eq:v60-joint-upper} instead.
Since $|q_w-L\mu_*|\le CL$, for $L\le\sqrt m$ its exponential
factor is at most $C e^{-c'|q-m\mu_*|^2/m}$. This follows from
$|v-d|^2\ge |v|^2/2-|d|^2$ and $m-L\ge m/2$.
Sum the outer cylinder cover, using $\pi_i\ge3/7$. It has weighted
length at most $C(|J\cap[0,1]|+\rho_0^L)$. This also controls
boundary atoms and proves \eqref{eq:v60-cylinder-upper}.
\end{proof}

\subsection{The original continuous density for every small perturbation}
Let $p_{m,\epsilon}^{f,d}$ denote the density of $S_m\tau_\epsilon$
with insertion $f(y_0)d(y_m)$, where $f,d$ are bounded Lipschitz
functions on $[0,1]$. Put $p_{m,\epsilon}=p_{m,\epsilon}^{1,1}$ and
$\|f\|_{\rm Lip}=\|f\|_\infty+\operatorname{Lip}(f)$.
For $r=t-3m$ and $\delta_w=r-A_m(w)-\epsilon B_m(w)$ the exact coarea
formula is
\begin{equation}\label{eq:v60-exact-perturbed-coarea}
 p_{m,\epsilon}^{f,d}(t)=
 \sum_w\frac{w_\mu}{1-a_w}
 f\left(\frac{b_w-\delta_w}{1-a_w}\right)
 d\left(b_w+a_w\frac{b_w-\delta_w}{1-a_w}\right)
 \one_{[b_w+a_w-1,b_w]}(\delta_w).
\end{equation}
The functions are evaluated only on the displayed support. The stable
derivative $1-a_w\ge1-\rho_0$ is independent of $\epsilon$ and the
word. Conditional uniformity of $y_0$, not independent roof noise,
proves this formula. In particular the unweighted integral is one.

Retain the one-periodic endpoint overlap $\Theta_{f,d}$ from
Theorem~\ref{thm:v59-markov-density} and define the convergent sum
\begin{equation}\label{eq:v60-arithmetic-profile}
 \mathcal R_{m,\epsilon}^{f,d}(t)=\frac{1}{\sqrt m}
 \sum_{l\in\Z}
 g_{\Sigma_*}\left(\frac{r-\epsilon l-3m/7}{\sqrt m},
                         \frac{l-2m/7}{\sqrt m}\right)
                         \Theta_{f,d}(r-\epsilon l).
\end{equation}
This profile retains the actual projected arithmetic for rational,
irrational and count-dependent perturbations.

\begin{theorem}[Uniform perturbed pointwise density]
\label{thm:v60-perturbed-density}
For every fixed $0<\epsilon_0<1/4$ and $m\ge2$,
\begin{align}
 \sup_{|\epsilon|\le\epsilon_0}\operatorname*{ess\,sup}_{t\in\R}
 \left|\sqrt m\,p_{m,\epsilon}^{f,d}(t)
                   -\mathcal R_{m,\epsilon}^{f,d}(t)\right|
 &\le C\|f\|_{\rm Lip}\|d\|_{\rm Lip}
                  \frac{\log(2+m)^{3/2}}{\sqrt m},
                                                   \label{eq:v60-weighted-pointwise}\\
 \sup_{|\epsilon|\le\epsilon_0}\operatorname*{ess\,sup}_{t\in\R}
 \left|\sqrt m\,p_{m,\epsilon}(t)
 -g_{\sigma_\epsilon^2}\left(\frac{t-m\bar\tau_\epsilon}{\sqrt m}\right)\right|
 &\le C\frac{\log(2+m)^{3/2}}{\sqrt m}.
                                                   \label{eq:v60-unweighted-pointwise}
\end{align}
Moreover $\sup_{m,|\epsilon|\le\epsilon_0}\sqrt m
\|p_{m,\epsilon}\|_\infty<\infty$. The endpoint-weighted arithmetic
profile is not replaced by a product of means.
\end{theorem}
\begin{proof}
The height upper bound follows at once from the exact coarea formula
and \eqref{eq:v60-strip-upper}, with $H=1$. For each fixed integer
pair $(k,l)$ set $\delta=r-k-\epsilon l$. Its unweighted coarea
contribution is trapped between
\[
 \mu\{Q_m=(k,l),b_m\in[\delta,\delta+1-\rho_0^m]\}
 \quad\hbox{and}\quad
 (1-\rho_0^m)^{-1}
 \mu\{Q_m=(k,l),b_m\in[\delta,\delta+1]\}.
\]
At a coarea point the two endpoints differ from $b_m-\delta,b_m$ by
$O(\rho_0^m)$. The weighted version follows by applying
\eqref{eq:v60-cylinder-LLT} to
$f(v-\delta)d(v)\one_{[\delta,\delta+1]}(v)$ through its bounded
variation. The supremum and variation of this function are bounded by
$C\|f\|_{\rm Lip}\|d\|_{\rm Lip}$, uniformly in $\delta$.
The changed interval boundaries are paid by
\eqref{eq:v60-cylinder-upper}. Thus, after multiplication by $m$,
the contribution equals
\[
 g_{\Sigma_*}((k-3m/7,l-2m/7)/\sqrt m)
 \int_{[0,1]\cap[\delta,\delta+1]} f(v-\delta)d(v)\,dv
 +O\left(\|f\|_{\rm Lip}\|d\|_{\rm Lip}
                \left(\frac{1+L}{\sqrt m}+\rho_0^L\right)\right).
\]
Choose $L=\lceil\log(m)/(2\log(4/3))\rceil$ for sufficiently large
$m$, so that $L\le\min(m/2,\sqrt m)$ and $\rho_0^L\le m^{-1/2}$.

It is important not to sum this error over $m$ possible values of $l$.
First restrict to $|l-2m/7|\le K\sqrt{m\log(2+m)}$.
For each $l$ at most three integers $k$ have $|\delta|\le1$.
The positive envelope \eqref{eq:v60-joint-upper} bounds the omitted
coarea contribution, after normalization by $\sqrt m$, by $Cm^{-2}$
when $K$ is fixed sufficiently large. The same estimate holds for
the omitted Gaussian sum. This follows by summing the Gaussian tail
in the second coordinate; it is uniform in $t$ and $\epsilon$.
There are $O(\sqrt{m\log(2+m)})$ retained pairs. Their total normalized
error is $O(\log(2+m)^{3/2}/\sqrt m)$.

On each retained pair replace $k$ in the Gaussian by $r-\epsilon l$.
The change of the normalized first coordinate is at most $m^{-1/2}$.
Its derivative is bounded by a summable Gaussian envelope in the second
coordinate, giving total error $O(m^{-1/2})$ after the same normalization.
Summing the integrals over $k$ gives exactly $\Theta_{f,d}(r-\epsilon l)$.
This proves \eqref{eq:v60-weighted-pointwise}. Complex weights use
variation and the positive envelope for domination, not an order
comparison of complex densities.

For $f=d=1$ the triangle partition identity gives $\Theta_{1,1}=1$.
Put $z=(t-m\bar\tau_\epsilon)/\sqrt m$ and
$y_l=(l-2m/7)/\sqrt m$. The remaining reference is the mesh
$m^{-1/2}$ Riemann sum of $F_{\epsilon,z}(y)=g_{\Sigma_*}(z-\epsilon y,y)$.
Uniformly in $z$ and $|\epsilon|\le\epsilon_0$,
\[
 \int_\R F_{\epsilon,z}(y)\,dy=g_{\sigma_\epsilon^2}(z),\qquad
 \int_\R |F'_{\epsilon,z}(y)|\,dy\le C.
\]
These follow by completing the Gaussian square; the conditional variance
is $\det\Sigma_*/\sigma_\epsilon^2$, bounded above and below.
The elementary interval-by-interval Riemann error is at most
$m^{-1/2}\int|F'|$. This proves \eqref{eq:v60-unweighted-pointwise}.
Finitely many small $m$ are absorbed using the uniform height bound.
No growing-band spectral constant or changing scalar lattice span
has entered the proof.
\end{proof}
